% ============================================================
%  TCC — Protocolo de Diagnóstico Estatístico para Séries Temporais
%  SENAI CIMATEC · Pós-Graduação em Ciência de Dados e IA
%  Vinícius da Fonseca Abreu · 2026
% ============================================================

\documentclass[
  12pt,
  openright,
  twoside,
  a4paper,
  english,
  brazil
]{abntex2}

% ---- Pacotes essenciais ----
\usepackage[T1]{fontenc}
\usepackage[utf8]{inputenc}
\usepackage{lmodern}
\usepackage{microtype}
\usepackage{indentfirst}
\usepackage{color}
\usepackage{graphicx}
\usepackage{booktabs}
\usepackage{amsmath}
\usepackage{amssymb}
\usepackage{hyperref}
\usepackage{listings}
\usepackage{float}
\usepackage{multirow}

% ---- Citações ABNT ----
\usepackage[brazilian,hyperpageref]{backref}
\usepackage[alf]{abntex2cite}

% ---- Configuração de margens ABNT ----
% Superior 3cm | Inferior 2cm | Esquerda 3cm | Direita 2cm
\usepackage[
  top=3cm,
  bottom=2cm,
  left=3cm,
  right=2cm
]{geometry}

% ---- Espaçamento ----
\setlength{\parindent}{1.25cm}
\setlength{\parskip}{0pt}
\OnehalfSpacing

% ---- Informações do trabalho ----
\titulo{Protocolo de Diagnóstico Estatístico para Seleção de Modelos em Séries Temporais: uma abordagem aplicada ao dataset M4}
\autor{Vinícius da Fonseca Abreu}
\local{Salvador}
\data{2026}
\orientador{Prof. Eldman}
\instituicao{%
  SENAI CIMATEC\\
  Pós-Graduação em Ciência de Dados e Inteligência Artificial
}
\tipotrabalho{Trabalho de Conclusão de Curso}
\preambulo{Trabalho de Conclusão de Curso apresentado ao Programa de Pós-Graduação em Ciência de Dados e Inteligência Artificial do SENAI CIMATEC como requisito parcial para obtenção do título de especialista.}

% ============================================================
%  DOCUMENTO
% ============================================================
\begin{document}

\selectlanguage{brazil}
\frenchspacing

% ---- Capa ----
\imprimircapa

% ---- Folha de rosto ----
\imprimirfolhaderosto*

% ---- Sumário ----
\tableofcontents*
\cleardoublepage

% ============================================================
%  CAPÍTULO 1 — INTRODUÇÃO
% ============================================================
\chapter{Introdução}
\label{cap:introducao}

A previsão de séries temporais ocupa um papel central em diversas áreas do conhecimento — da economia à engenharia, passando pela saúde pública e pelo mercado financeiro. A capacidade de antecipar o comportamento futuro de uma variável a partir de seu histórico é, ao mesmo tempo, um problema matemático bem definido e uma demanda prática recorrente. Conforme apontado por \citeonline{makridakis2020m4}, o crescimento expressivo do número de métodos disponíveis — incluindo técnicas de aprendizado de máquina e aprendizado profundo — ampliou as opções do pesquisador sem, contudo, oferecer critérios claros para a escolha entre elas. Essa lacuna metodológica é o ponto de partida deste trabalho: \textit{como escolher qual modelo usar?}

A prática corrente em muitos estudos e projetos aplicados é operar por tentativa e erro — aplica-se ARIMA, testa-se LSTM, comparam-se os erros de validação, e o modelo com menor métrica é declarado vencedor. Essa abordagem ignora uma etapa anterior fundamental: entender a natureza estatística da série antes de qualquer modelagem. Tratar uma série não-linear com ARIMA, que pressupõe linearidade, ou aplicar redes neurais a séries estacionárias simples, onde modelos clássicos generalizam melhor, são erros metodológicos que não aparecem nos resultados de treino, mas se manifestam na previsão real. Estudo empírico de \citeonline{makridakis2018statistical} sobre um grande conjunto de séries reais evidenciou que métodos estatísticos clássicos superam consistentemente técnicas de aprendizado de máquina em horizontes curtos e séries com padrões estáveis, enquanto o contrário ocorre em séries com forte não-linearidade — resultado que reforça a necessidade de diagnóstico estatístico prévio à escolha do modelo.

Este trabalho parte de uma premissa diferente: \textbf{primeiro diagnosticar a série, depois escolher o modelo}. A hipótese central é que a seleção de modelo guiada por um protocolo de diagnóstico estatístico — composto por detecção de sazonalidade (Etapa 0) e pelos testes ADF, KPSS, BDS, ARCH e pelo expoente de Hurst — produz previsões mais acuradas do que o benchmark automático (Auto-ARIMA aplicado cegamente), medido pelas métricas oficiais da Competição M4: sMAPE, MASE e OWA.

O protocolo proposto funciona como um funil de cinco etapas aplicado a cada série temporal individualmente, antes da fase de modelagem. A etapa zero detecta sazonalidade pela força sazonal da decomposição STL, produzindo um indicador binário que determina se o modelo final deve incluir componentes sazonais. A primeira etapa avalia a estacionariedade por meio dos testes ADF e KPSS em conjunto, evitando falsos positivos que ocorrem quando apenas um dos testes é utilizado. A segunda etapa aplica o teste BDS nos resíduos de um modelo autorregressivo ajustado, identificando estrutura não-linear — condição que invalida o uso de modelos lineares. A terceira etapa investiga heterocedasticidade condicional via teste ARCH de Engle, cujo resultado direciona para modelos da família GARCH. A quarta etapa estima a memória longa da série pelo expoente de Hurst combinado com o estimador GPH, indicando o uso de ARFIMA quando pertinente.

O experimento é conduzido sobre o subconjunto mensal da Competição M4, que reúne aproximadamente 48.000 séries temporais univariadas reais de domínios variados. A M4 é o maior benchmark público de previsão de séries temporais disponível na literatura \cite{makridakis2020m4}, e seu uso garante comparabilidade com dezenas de trabalhos anteriores que utilizaram o mesmo conjunto de dados.

\section{Motivação}
\label{sec:motivacao}

A escolha de modelos por critério estatístico, em vez de por validação cruzada pura, não é nova. Já em \citeonline{box1970time}, a identificação da ordem do modelo ARIMA era feita por inspeção das funções de autocorrelação e autocorrelação parcial — um diagnóstico visual que antecede a estimação. O que mudou nas últimas décadas foi a disponibilidade de modelos alternativos: redes neurais recorrentes, arquiteturas de atenção e modelos baseados em gradiente boosting expandiram o espaço de escolha de forma substancial, sem que a metodologia de seleção tenha acompanhado esse crescimento na mesma proporção.

A Competição M4 revelou um resultado que ainda divide opiniões na comunidade: métodos híbridos e métodos clássicos bem calibrados superaram, em média, as redes neurais profundas aplicadas sem restrições \cite{makridakis2018m4results}. Esse achado sugere que a qualidade do modelo não depende apenas da sua capacidade expressiva, mas do alinhamento entre as propriedades do modelo e as propriedades estatísticas da série que ele tenta capturar. Nessa mesma competição, o método FFORMA de \citeonline{montero2020fforma} — que seleciona e pondera modelos automaticamente com base em características estatísticas extraídas de cada série — ficou em segundo lugar geral, validando empiricamente a hipótese de que a caracterização prévia da série é fator determinante para a acurácia da previsão.

O teste BDS \cite{broock1996test}, desenvolvido originalmente para detectar dependência não-linear em resíduos de modelos econométricos, oferece um critério formal para decidir quando um modelo linear é suficiente e quando a estrutura da série exige abordagens não-paramétricas. A ausência desse tipo de triagem prévia na maioria dos pipelines de previsão automatizados representa uma lacuna que este trabalho busca endereçar.

\section{Objetivo Geral}
\label{sec:objetivo-geral}

Propor e validar um protocolo de diagnóstico estatístico que guie a seleção de modelos de previsão para séries temporais, comparando seu desempenho frente ao benchmark automático (Auto-ARIMA) sobre o subconjunto mensal da Competição M4.

\section{Objetivos Específicos}
\label{sec:objetivos-especificos}

\begin{enumerate}
  \item Implementar o pipeline de diagnóstico composto pela detecção de sazonalidade (Etapa 0) e pelos testes ADF, KPSS, BDS, ARCH e expoente de Hurst;
  \item Classificar as séries do M4 mensal segundo os resultados do diagnóstico e atribuir o modelo indicado para cada perfil;
  \item Treinar e avaliar os modelos atribuídos — ARIMA, SARIMA, LSTM, ARIMA-GARCH e ARFIMA — sobre as séries correspondentes;
  \item Comparar os resultados em sMAPE, MASE e OWA frente a dois benchmarks cegos: Auto-ARIMA e Método Theta aplicados a todas as séries sem diagnóstico prévio;
  \item Avaliar a acurácia diagnóstica do protocolo, medindo a frequência com que o modelo indicado corresponde ao modelo de melhor desempenho empírico para cada série;
  \item Discutir os trade-offs de custo computacional, interpretabilidade e acurácia entre as abordagens comparadas.
\end{enumerate}

\section{Estrutura do Trabalho}
\label{sec:estrutura}

O trabalho está organizado em cinco capítulos. O Capítulo \ref{cap:fundamentacao} apresenta a fundamentação teórica, cobrindo os conceitos de séries temporais, os testes estatísticos utilizados no diagnóstico, os modelos de previsão abordados e as métricas de avaliação. O Capítulo 3 descreve a metodologia, detalhando o protocolo de diagnóstico e o design experimental. O Capítulo 4 apresenta e discute os resultados. O Capítulo 5 traz as conclusões e aponta direções para trabalhos futuros.

% ============================================================
%  CAPÍTULO 2 — FUNDAMENTAÇÃO TEÓRICA
% ============================================================
\chapter{Fundamentação Teórica}
\label{cap:fundamentacao}

\section{Séries Temporais}
\label{sec:series-temporais}

Segundo \citeonline{morettin2006analise}, uma série temporal é uma sequência de observações de uma variável aleatória indexadas no tempo, $\{y_t\}_{t=1}^{T}$, em que a ordem das observações carrega informação. Essa dependência temporal é o que diferencia séries temporais de amostras aleatórias independentes e identicamente distribuídas, tornando sua modelagem metodologicamente distinta da estatística clássica.

Conforme \citeonline{hyndman2018forecasting}, uma série temporal pode ser decomposta em quatro componentes: tendência ($T_t$), sazonalidade ($S_t$), ciclo ($C_t$) e componente irregular ou ruído ($\varepsilon_t$). No modelo aditivo, esses componentes se somam:

\begin{equation}
  y_t = T_t + S_t + C_t + \varepsilon_t
  \label{eq:decomposicao-aditiva}
\end{equation}

No modelo multiplicativo, interagem de forma proporcional:

\begin{equation}
  y_t = T_t \cdot S_t \cdot C_t \cdot \varepsilon_t
  \label{eq:decomposicao-multiplicativa}
\end{equation}

De acordo com \citeonline{hyndman2018forecasting}, a escolha entre os dois depende da estrutura dos dados: quando a amplitude da sazonalidade cresce proporcionalmente à tendência, o modelo multiplicativo é mais adequado.

\subsection{Processos Estocásticos e Estacionariedade}
\label{subsec:estacionariedade}

Conforme \citeonline{hamilton1994time}, um processo estocástico $\{y_t\}$ é dito \textbf{estritamente estacionário} se a distribuição conjunta de qualquer subconjunto de variáveis é invariante a deslocamentos no tempo. Na prática, trabalha-se com a condição mais fraca de \textbf{estacionariedade fraca} (ou de segunda ordem), que exige \cite{hamilton1994time}:

\begin{enumerate}
  \item $\mathbb{E}[y_t] = \mu$, constante para todo $t$;
  \item $\text{Var}(y_t) = \sigma^2 < \infty$, constante para todo $t$;
  \item $\text{Cov}(y_t, y_{t-k}) = \gamma(k)$, dependente apenas da defasagem $k$, não de $t$.
\end{enumerate}

A estacionariedade é um pressuposto fundamental dos modelos ARIMA e ARMA. Quando uma série apresenta raiz unitária — isto é, quando seus choques têm efeitos permanentes — ela é integrada de ordem $d$, denotada $I(d)$, e precisa ser diferenciada $d$ vezes para tornar-se estacionária \cite{hamilton1994time}.

\section{Testes de Diagnóstico Estatístico}
\label{sec:testes-diagnostico}

\subsection{Teste de Dickey-Fuller Aumentado (ADF)}
\label{subsec:adf}

O Teste de Dickey-Fuller Aumentado (ADF) tem como base o teste original de \citeonline{dickey1979distribution}, que verificava a presença de raiz unitária em processos autorregressivos de primeira ordem. A extensão para a forma aumentada — aplicável a modelos ARMA de ordem desconhecida — foi desenvolvida por \citeonline{said1984testing}. A hipótese nula ($H_0$) postula que a série possui raiz unitária — ou seja, é não-estacionária — enquanto a hipótese alternativa ($H_1$) é a estacionariedade.

O modelo estimado na forma aumentada, conforme \citeonline{said1984testing}, é:

\begin{equation}
  \Delta y_t = \alpha + \beta t + \gamma y_{t-1} + \sum_{i=1}^{p} \delta_i \Delta y_{t-i} + \varepsilon_t
  \label{eq:adf}
\end{equation}

onde $\Delta y_t = y_t - y_{t-1}$ é o operador de primeira diferença, $\alpha$ é a constante, $\beta t$ é a tendência determinística, e os termos $\Delta y_{t-i}$ são incluídos para corrigir autocorrelação nos resíduos. O teste é sobre $H_0: \gamma = 0$; rejeitar $H_0$ implica que a série é estacionária.

A quantidade de defasagens $p$ é tipicamente selecionada via Critério de Informação de Akaike (AIC) ou Critério de Schwarz (BIC) \cite{said1984testing}.

\subsection{Teste KPSS}
\label{subsec:kpss}

O teste de Kwiatkowski-Phillips-Schmidt-Shin, proposto por \citeonline{kwiatkowski1992testing}, complementa o ADF ao inverter as hipóteses: $H_0$ é que a série é estacionária (ou trend-estacionária) e $H_1$ é que ela possui raiz unitária. Conforme os autores, a estatística de teste é construída a partir dos resíduos de uma regressão de $y_t$ sobre uma constante ou sobre uma tendência linear \cite{kwiatkowski1992testing}.

Conforme \citeonline{kwiatkowski1992testing}, o KPSS foi desenvolvido precisamente como complemento ao ADF: ao inverter as hipóteses, os dois testes se controlam mutuamente. Quando concordam — ADF rejeita $H_0$ e KPSS não rejeita $H_0$ —, a evidência de estacionariedade é robusta. Quando divergem, a série pode apresentar raiz unitária fraca ou quebra estrutural, o que requer investigação adicional \cite{kwiatkowski1992testing}.

\subsection{Teste BDS}
\label{subsec:bds}

O teste de Brock-Dechert-Scheinkman, desenvolvido por \citeonline{broock1996test}, é um teste não-paramétrico de independência serial que detecta estrutura não-linear nos dados. Aplicado sobre os resíduos de um modelo linear ajustado — tipicamente um AR —, sua hipótese nula postula que os resíduos são independentes e identicamente distribuídos (i.i.d.). Rejeitar $H_0$ indica que o modelo linear não capturou toda a estrutura da série, sugerindo dependência não-linear \cite{broock1996test}.

Conforme \citeonline{broock1996test}, a estatística BDS é fundamentada na \textit{integral de correlação}, definida como a fração de pares de vetores de dimensão $m$ que estão a uma distância menor que $\varepsilon$ no espaço de estados:

\begin{equation}
  C_{m,T}(\varepsilon) = \frac{2}{T_m(T_m - 1)} \sum_{t < s} \mathbf{1}\left(\|y_t^m - y_s^m\| < \varepsilon\right)
  \label{eq:integral-correlacao}
\end{equation}

onde $T_m = T - m + 1$, $y_t^m = (y_t, y_{t+1}, \ldots, y_{t+m-1})$ é o vetor de imersão de dimensão $m$ e $\mathbf{1}(\cdot)$ é a função indicadora. Sob a hipótese nula de i.i.d., demonstra-se que $C_{m,T}(\varepsilon) \to [C_{1,T}(\varepsilon)]^m$, e a distância entre esses valores, devidamente normalizada, converge para uma distribuição normal padrão \cite{broock1996test}.

Neste trabalho, a rejeição do BDS sobre os resíduos de um AR ajustado é interpretada como evidência de que a série possui estrutura não-linear, invalidando o ARIMA como modelo final e direcionando para o modelo LSTM.

\subsection{Teste ARCH de Engle}
\label{subsec:arch}

A heterocedasticidade condicional é uma propriedade comum em séries financeiras e econômicas, na qual a variância dos erros não é constante ao longo do tempo, mas depende de erros passados. Conforme \citeonline{engle1982autoregressive}, o teste ARCH verifica se os resíduos ao quadrado de um modelo ajustado apresentam autocorrelação, o que constitui evidência de volatilidade condicional.

O modelo ARCH($q$), proposto por \citeonline{engle1982autoregressive}, expressa a variância condicional como:

\begin{equation}
  \sigma_t^2 = \alpha_0 + \sum_{i=1}^{q} \alpha_i \varepsilon_{t-i}^2
  \label{eq:arch}
\end{equation}

A hipótese nula do teste é $H_0: \alpha_1 = \alpha_2 = \cdots = \alpha_q = 0$, isto é, ausência de efeito ARCH. A detecção de heterocedasticidade condicional direciona para a família de modelos GARCH, que \citeonline{bollerslev1986generalized} generalizou ao permitir que a variância condicional dependa também de suas próprias defasagens passadas.

\subsection{Expoente de Hurst e Memória Longa}
\label{subsec:hurst}

Séries com \textbf{memória longa} (ou dependência de longo alcance) apresentam autocorrelações que decaem em lei de potência — muito mais lentamente do que o decaimento exponencial observado em processos ARMA. Essa propriedade foi formalizada por \citeonline{hurst1951long-storage} por meio do expoente $H$, estimado pela análise R/S (\textit{Rescaled Range}) ou pelo estimador GPH de \citeonline{geweke1983estimation}, que opera no domínio da frequência via regressão log-periódograma.

De acordo com \citeonline{hurst1951long-storage} e \citeonline{granger1980introduction}, os valores de $H$ se interpretam da seguinte forma:

\begin{itemize}
  \item $H = 0{,}5$: processo sem memória, equivalente ao ruído branco;
  \item $H > 0{,}5$: memória longa positiva — valores altos tendem a ser seguidos de valores altos;
  \item $H < 0{,}5$: antipersistência — reversão à média mais rápida que o esperado.
\end{itemize}

Conforme \citeonline{granger1980introduction}, séries com $H > 0{,}5$ possuem dependência de longo prazo cuja intensidade justifica o uso de modelos fracionalmente integrados, como o ARFIMA, em substituição ao ARIMA padrão, que opera com integrações de ordem inteira e não captura adequadamente essa estrutura.

\section{Modelos de Previsão}
\label{sec:modelos-previsao}

\subsection{ARIMA}
\label{subsec:arima}

O modelo Autorregressivo Integrado de Médias Móveis, ARIMA($p, d, q$), foi proposto por \citeonline{box1970time} e constitui o modelo clássico de referência para séries temporais univariadas lineares. Conforme os autores, o modelo combina três componentes:

\begin{itemize}
  \item \textbf{AR($p$)}: a observação atual é uma combinação linear de $p$ observações passadas;
  \item \textbf{I($d$)}: a série é diferenciada $d$ vezes até atingir estacionariedade;
  \item \textbf{MA($q$)}: a observação atual depende de $q$ erros de previsão passados.
\end{itemize}

A representação geral do modelo, após $d$ diferenciações, é \cite{box1970time}:

\begin{equation}
  \phi(B)(1 - B)^d y_t = \theta(B)\varepsilon_t
  \label{eq:arima}
\end{equation}

onde $B$ é o operador de retardo ($By_t = y_{t-1}$), $\phi(B) = 1 - \phi_1 B - \cdots - \phi_p B^p$ e $\theta(B) = 1 + \theta_1 B + \cdots + \theta_q B^q$. Neste trabalho, o Auto-ARIMA — que automatiza a seleção de $p, d, q$ via minimização do AIC ou BIC — é utilizado como benchmark de comparação.

A extensão sazonal, SARIMA($p, d, q$)($P, D, Q$)$_s$, incorpora componentes análogos operando na frequência sazonal $s$, conforme \citeonline{box1970time}.

\subsection{ARFIMA}
\label{subsec:arfima}

O modelo ARFIMA (Autorregressivo Fracionalmente Integrado de Médias Móveis), introduzido por \citeonline{granger1980introduction}, generaliza o ARIMA ao permitir que o parâmetro de integração $d$ assuma valores não inteiros $d \in (-0{,}5;\, 0{,}5)$:

\begin{equation}
  \phi(B)(1 - B)^d y_t = \theta(B)\varepsilon_t, \quad d \in \mathbb{R}
  \label{eq:arfima}
\end{equation}

Conforme \citeonline{granger1980introduction}, quando $0 < d < 0{,}5$, o processo é estacionário e apresenta memória longa, com o expoente de Hurst relacionando-se a $d$ pela expressão $H = d + 0{,}5$. Neste trabalho, o ARFIMA é indicado pelo protocolo quando o diagnóstico da Etapa 4 identifica $H > 0{,}7$, limiar adotado como critério operacional para memória longa pronunciada.

\subsection{Modelos GARCH}
\label{subsec:garch}

O modelo GARCH($p, q$), proposto por \citeonline{bollerslev1986generalized}, modela a variância condicional como uma função das variâncias e erros ao quadrado passados:

\begin{equation}
  \sigma_t^2 = \omega + \sum_{i=1}^{q} \alpha_i \varepsilon_{t-i}^2 + \sum_{j=1}^{p} \beta_j \sigma_{t-j}^2
  \label{eq:garch}
\end{equation}

A especificação GARCH(1,1) — caso particular com $p = q = 1$ — tem se mostrado empiricamente competitiva frente a especificações mais complexas em uma ampla variedade de séries financeiras, conforme demonstrado por \citeonline{hansen2005forecast} em estudo comparativo de modelos de volatilidade. Neste trabalho, o GARCH é combinado com ARIMA — o ARIMA modela a média condicional e o GARCH modela os resíduos — quando a Etapa 3 do protocolo detecta heterocedasticidade condicional.

\subsection{LSTM — Long Short-Term Memory}
\label{subsec:lstm}

As redes LSTM, propostas por \citeonline{hochreiter1997long}, são uma variante das redes neurais recorrentes (RNNs) desenvolvida para capturar dependências de longo prazo em sequências, contornando o problema de dissipação do gradiente que afeta RNNs simples. Na formulação original, a arquitetura contava com portões de entrada e saída; o \textit{forget gate} — que permite à célula redefinir seu estado interno — foi introduzido posteriormente por \citeonline{gers2000learning}. Na formulação padrão atual, com os três portões, a unidade LSTM opera da seguinte forma \cite{hochreiter1997long,gers2000learning}:

\begin{align}
  f_t &= \sigma(W_f [h_{t-1}, x_t] + b_f) \label{eq:lstm-forget} \\
  i_t &= \sigma(W_i [h_{t-1}, x_t] + b_i) \label{eq:lstm-input} \\
  \tilde{C}_t &= \tanh(W_C [h_{t-1}, x_t] + b_C) \label{eq:lstm-cell-candidate} \\
  C_t &= f_t \odot C_{t-1} + i_t \odot \tilde{C}_t \label{eq:lstm-cell} \\
  o_t &= \sigma(W_o [h_{t-1}, x_t] + b_o) \label{eq:lstm-output} \\
  h_t &= o_t \odot \tanh(C_t) \label{eq:lstm-hidden}
\end{align}

onde $\sigma$ é a função sigmoide, $\odot$ é o produto de Hadamard (elemento a elemento), $x_t$ é a entrada no instante $t$ e $h_t$ é o estado oculto transmitido para o passo seguinte \cite{hochreiter1997long}.

A ausência de pressupostos de linearidade ou distribuição dos dados torna o LSTM adequado para séries onde o teste BDS rejeita a hipótese de i.i.d. nos resíduos, indicando estrutura não-linear. A Competição M4 mostrou que modelos híbridos combinando Suavização Exponencial com LSTM — o ES-RNN de \citeonline{smyl2020hybrid}, vencedor da competição — superaram tanto modelos puramente clássicos quanto puramente neurais. Revisão abrangente da literatura por \citeonline{hewamalage2021recurrent} indica que, em previsão de séries univariadas de curto horizonte, LSTMs raramente superam métodos estatísticos bem especificados de forma consistente, sendo mais eficazes em séries longas com padrões complexos e não-lineares — o que reforça a lógica de usar o BDS como critério de triagem antes de optar por arquiteturas recorrentes.

\subsection{Método Theta}
\label{subsec:theta}

O Método Theta, proposto por \citeonline{assimakopoulos2000theta}, é um método de decomposição que modifica a curvatura local de uma série temporal por meio de um coeficiente $\theta$ aplicado às segundas diferenças. A operação define as \textit{theta lines}, curvas que satisfazem:

\begin{equation}
  \Delta^2 \Theta_\theta(t) = \theta \cdot \Delta^2 y_t
  \label{eq:theta-def}
\end{equation}

Na versão padrão do método, utilizam-se dois valores complementares. Para $\theta = 0$, a theta line é a reta de regressão linear ajustada a $y_t$, denotada $\Theta_0(t) = a + bt$, que captura a tendência de longo prazo. Para $\theta = 2$, obtém-se $\Theta_2(t) = 2y_t - \Theta_0(t)$, série com curvatura local duplicada que preserva as oscilações de curto prazo \cite{assimakopoulos2000theta}. A previsão final é a média aritmética das extrapolações independentes das duas theta lines:

\begin{equation}
  \hat{y}_{T+h} = \frac{1}{2}\,\hat{\Theta}_0(T+h) + \frac{1}{2}\,\hat{\Theta}_2(T+h)
  \label{eq:theta-forecast}
\end{equation}

onde $\Theta_0$ é extrapolada pelo crescimento linear e $\Theta_2$ é extrapolada por Suavização Exponencial Simples (SES).

Conforme demonstrado por \citeonline{hyndman2003unmasking}, a previsão resultante é algebricamente equivalente a SES com \textit{drift}, isto é, SES com acréscimo de uma tendência linear:

\begin{equation}
  \hat{y}_{T+h} = \hat{y}_{T}^{\text{SES}} + \frac{\hat{b}}{2} \cdot h
  \label{eq:theta-ses-drift}
\end{equation}

onde $\hat{b}$ é o coeficiente angular da regressão linear ajustada a $y_t$ e $\hat{y}_{T}^{\text{SES}}$ é o nível estimado pelo SES aplicado a $\Theta_2$ \cite{hyndman2003unmasking}. Para séries sazonais, a sazonalidade é removida antes da aplicação do método e reintroduzida após a previsão.

O Método Theta venceu a Competição M3 \cite{assimakopoulos2000theta} e manteve desempenho competitivo na M4 \cite{makridakis2018m4results}, o que o torna referência obrigatória como benchmark em estudos sobre o subconjunto mensal utilizado neste trabalho. Sua principal vantagem é a ausência de hiperparâmetros de estrutura: o único parâmetro livre é o coeficiente de suavização do SES, estimado por mínimos quadrados \cite{hyndman2003unmasking}.

\section{A Competição M4}
\label{sec:m4}

Conforme \citeonline{makridakis2020m4}, a Competição M4 é a quarta edição da série de competições de previsão de séries temporais iniciada por Makridakis em 1982. A edição de 2018 reuniu 100.000 séries temporais univariadas reais de seis frequências temporais (anual, trimestral, mensal, semanal, diária e horária) e 61 métodos competidores, constituindo o maior benchmark público de previsão disponível na literatura \cite{makridakis2020m4}.

O subconjunto mensal, utilizado neste trabalho, contém aproximadamente 48.000 séries de domínios variados — economia, finanças, indústria, demografia e outros —, com horizonte de previsão de 18 passos à frente \cite{makridakis2020m4}.

De acordo com \citeonline{makridakis2018m4results}, os principais achados da competição relevantes para este trabalho são:

\begin{enumerate}
  \item Métodos híbridos — combinando modelos estatísticos e redes neurais — foram os mais acurados no geral;
  \item A combinação de previsões (\textit{forecast combination}) produziu resultados robustos mesmo quando os métodos individuais eram relativamente simples;
  \item Modelos clássicos bem calibrados superaram redes neurais aplicadas sem restrições em diversas categorias;
  \item A identificação de características da série antes da modelagem apareceu como fator diferencial nos métodos mais bem posicionados.
\end{enumerate}

\section{Métricas de Avaliação}
\label{sec:metricas}

A avaliação dos modelos segue as métricas oficiais da Competição M4.

\subsection{sMAPE — Symmetric Mean Absolute Percentage Error}
\label{subsec:smape}

Proposto por \citeonline{makridakis1993accuracy}, o sMAPE corrige a assimetria do MAPE tradicional ao dividir o erro pelo semissoma dos valores real e previsto:

\begin{equation}
  \text{sMAPE} = \frac{100\%}{h} \sum_{t=1}^{h} \frac{|y_t - \hat{y}_t|}{(|y_t| + |\hat{y}_t|)/2}
  \label{eq:smape}
\end{equation}

onde $h$ é o horizonte de previsão e $\hat{y}_t$ é o valor previsto \cite{makridakis1993accuracy}. Valores menores indicam maior acurácia. A métrica é adimensional e comparável entre séries de escalas diferentes.

\subsection{MASE — Mean Absolute Scaled Error}
\label{subsec:mase}

Proposto por \citeonline{hyndman2006another}, o MASE escala o erro absoluto médio pelo erro de um método de referência ingênuo (naive) calculado na amostra de treino:

\begin{equation}
  \text{MASE} = \frac{1}{h} \sum_{t=1}^{h} \frac{|y_t - \hat{y}_t|}{\frac{1}{T - m} \sum_{t=m+1}^{T} |y_t - y_{t-m}|}
  \label{eq:mase}
\end{equation}

onde $m$ é o período sazonal — 12 para séries mensais \cite{hyndman2006another}. MASE $< 1$ indica que o modelo supera o naive sazonal. Conforme \citeonline{hyndman2006another}, ao contrário do MAPE, o MASE é bem definido mesmo quando $y_t = 0$, o que o torna mais adequado para séries com valores nulos ou próximos de zero.

\subsection{OWA — Overall Weighted Average}
\label{subsec:owa}

O OWA é a métrica síntese adotada pela Competição M4 \cite{makridakis2018m4results}. Conforme \citeonline{makridakis2018m4results}, ela é calculada como a média das razões entre o desempenho do método avaliado e o desempenho de um benchmark de referência (Naive 2):

\begin{equation}
  \text{OWA} = \frac{1}{2}\left(\frac{\text{sMAPE}}{\text{sMAPE}_{\text{Naive2}}} + \frac{\text{MASE}}{\text{MASE}_{\text{Naive2}}}\right)
  \label{eq:owa}
\end{equation}

OWA $< 1$ significa que o método supera o benchmark \cite{makridakis2018m4results}. É a métrica primária utilizada neste trabalho para comparar o protocolo guiado com o Auto-ARIMA.

% ============================================================
%  CAPÍTULO 3 — METODOLOGIA
% ============================================================
\chapter{Metodologia}
\label{cap:metodologia}

Este capítulo descreve o delineamento da pesquisa, o conjunto de dados utilizado, o protocolo de diagnóstico estatístico proposto, a família de modelos candidatos, o baseline de comparação e o pipeline experimental adotado para validar a hipótese central do trabalho.

% ----------------------------------------------------------
\section{Delineamento da Pesquisa}
\label{sec:delineamento}

A pesquisa tem caráter quantitativo e natureza aplicada, com delineamento experimental comparativo \cite{gil2002pesquisa}. A estratégia consiste em aplicar dois pipelines de seleção de modelos ao mesmo conjunto de séries temporais e comparar seus desempenhos pelas métricas oficiais da Competição M4.

O pipeline proposto — denominado aqui \textit{pipeline guiado} — aplica o protocolo de diagnóstico estatístico de cinco etapas descrito na Seção \ref{sec:protocolo} para determinar, a priori, qual classe de modelo é mais adequada a cada série. O pipeline de referência — denominado \textit{baseline cego} — aplica Auto-ARIMA a todas as séries sem diagnóstico prévio. A comparação entre os dois pipelines constitui o teste empírico da hipótese enunciada na Seção \ref{sec:objetivo-geral}. A abordagem é conceitualmente próxima ao FFORMA \cite{montero2020fforma}, que utiliza características estatísticas das séries para ponderar modelos de previsão, diferindo por adotar testes de hipóteses formais — em vez de aprendizado supervisionado — como mecanismo de seleção, o que garante maior interpretabilidade e rastreabilidade das decisões do protocolo.

A unidade de análise é a série temporal individual. Para cada série são registrados: o perfil diagnóstico atribuído pelo protocolo, o modelo indicado, e os valores de sMAPE, MASE e OWA apurados no horizonte de 18 passos à frente. Os resultados são agregados globalmente e por perfil diagnóstico para identificar em quais classes de séries a seleção guiada apresenta ganho mais expressivo.

% ----------------------------------------------------------
\section{Conjunto de Dados}
\label{sec:dados}

O experimento é conduzido sobre o subconjunto mensal da Competição M4, que reúne 48.000 séries temporais univariadas reais de domínios variados — economia, finanças, indústria, demografia, entre outros \cite{makridakis2020m4}. As séries são disponibilizadas com divisão pré-definida entre amostra de treino e amostra de teste, e o horizonte de previsão é de $h = 18$ passos à frente \cite{makridakis2020m4}. A Tabela \ref{tab:m4-mensal} sintetiza as principais características do subconjunto.

\begin{table}[H]
  \centering
  \caption{Características do subconjunto mensal da Competição M4.}
  \label{tab:m4-mensal}
  \begin{tabular}{lc}
    \toprule
    \textbf{Atributo} & \textbf{Valor} \\
    \midrule
    Número de séries             & 48.000 \\
    Frequência                   & Mensal \\
    Horizonte de previsão        & 18 passos \\
    Comprimento mínimo (treino)  & 42 observações \\
    Comprimento máximo (treino)  & 2.794 observações \\
    Domínios                     & Macro, micro, finanças, indústria, demografia, outros \\
    \bottomrule
  \end{tabular}
  \fonte{\citeonline{makridakis2020m4}}
\end{table}

Não são realizadas transformações globais no conjunto. Cada série é tratada individualmente no pipeline de diagnóstico, e transformações locais — como diferenciação ou logaritmização — são aplicadas apenas quando indicadas pelos testes estatísticos em cada etapa.

O dataset inclui o rótulo de domínio de origem para cada série (Macro, Micro, Finanças, Indústria, Demografia, Outros), além do comprimento da série de treino, que varia entre 42 e 2.794 observações. Esses atributos são preservados ao longo do experimento para possibilitar desagregações analíticas descritas na Seção \ref{sec:analises-complementares}, onde o desempenho do protocolo é examinado separadamente por domínio e faixa de comprimento.

% ----------------------------------------------------------
\section{Protocolo de Diagnóstico Estatístico}
\label{sec:protocolo}

O protocolo proposto opera como um funil de cinco etapas sequenciais, aplicado a cada série individualmente antes de qualquer modelagem. O resultado de cada etapa determina se a série avança para a etapa seguinte ou é direcionada para um modelo específico.

\subsection{Etapa 0 — Sazonalidade}
\label{subsec:etapa0}

A etapa zero precede o funil de diagnóstico e determina se a série apresenta padrão sazonal relevante. A detecção é realizada pela decomposição STL (\textit{Seasonal and Trend decomposition using Loess}) com período $s = 12$, que separa a série nos componentes de tendência, sazonalidade e resíduo. A partir dessa decomposição, calcula-se a \textbf{força sazonal} $F_S$, conforme \citeonline{hyndman2018forecasting}:

\begin{equation}
  F_S = \max\!\left(0,\; 1 - \frac{\text{Var}(R_t)}{\text{Var}(S_t + R_t)}\right)
  \label{eq:forca-sazonal}
\end{equation}

onde $S_t$ é o componente sazonal e $R_t$ é o resíduo da decomposição. Quando $F_S > 0{,}64$ — limiar proposto por \citeonline{hyndman2018forecasting} para classificar uma série como ``sazonalmente forte'' —, a série recebe o indicador $\texttt{sazonal} = \text{verdadeiro}$, que é propagado ao final do protocolo: se nenhuma das Etapas 1–4 redirecionar a série para outro perfil, o modelo indicado é SARIMA em vez de ARIMA. Séries com $F_S \leq 0{,}64$ recebem $\texttt{sazonal} = \text{falso}$ e prosseguem normalmente pelo funil.

\subsection{Etapa 1 — Estacionariedade}
\label{subsec:etapa1}

A primeira etapa avalia se a série é estacionária em covariância por meio da aplicação conjunta dos testes ADF e KPSS. Conforme \citeonline{kwiatkowski1992testing}, o teste KPSS foi desenvolvido precisamente para complementar o ADF ao inverter as hipóteses: enquanto o ADF testa a raiz unitária como hipótese nula, o KPSS postula a estacionariedade como hipótese nula. A concordância entre os dois testes — ADF rejeita $H_0$ e KPSS não rejeita $H_0$ — constitui evidência robusta de estacionariedade; a discordância indica a presença de raiz unitária fraca ou quebra estrutural, e requer investigação adicional \cite{kwiatkowski1992testing}.

Quando a série não é estacionária, aplica-se diferenciação de primeira ordem e os testes são reaplicados. O número de diferenciações necessárias determina o parâmetro $d$ dos modelos ARIMA e ARFIMA. O nível de significância adotado é $\alpha = 0{,}05$ em ambos os testes.

\subsection{Etapa 2 — Linearidade}
\label{subsec:etapa2}

A segunda etapa investiga se a dinâmica da série pode ser capturada por um modelo linear. Um modelo autorregressivo de ordem $p$ — selecionado por AIC com $p_{\max} = 12$, compatível com a periodicidade mensal — é ajustado à série estacionarizada. O teste BDS de \citeonline{broock1996test} é então aplicado sobre os resíduos desse modelo para verificar a hipótese de independência serial.

De acordo com \citeonline{broock1996test}, a escolha do parâmetro de distância $\varepsilon$ e da dimensão de imersão $m$ influencia a potência do teste. Adota-se a faixa recomendada pelos autores: $m \in \{2, 3, 4, 5\}$ e $\varepsilon \in \{0{,}5;\, 1{,}0;\, 1{,}5\}$ vezes o desvio-padrão dos resíduos.

A rejeição consistente de $H_0$ em múltiplas combinações $(m, \varepsilon)$, ao nível $\alpha = 0{,}05$, é interpretada como evidência de que o modelo linear não captura toda a estrutura da série, indicando dependência não-linear. Nesse caso, a série é direcionada para modelos da família não-linear (LSTM). Quando $H_0$ não é rejeitada de forma consistente, a série avança para a Etapa 3.

\subsection{Etapa 3 — Heterocedasticidade Condicional}
\label{subsec:etapa3}

A terceira etapa verifica a presença de volatilidade condicional nos resíduos do modelo linear. Conforme \citeonline{engle1982autoregressive}, o teste ARCH examina se os resíduos ao quadrado apresentam autocorrelação significativa, evidenciando que a variância dos erros varia ao longo do tempo. O teste é aplicado com $q = 12$ defasagens sobre os resíduos do AR ajustado na Etapa 2, e o nível de significância adotado é $\alpha = 0{,}05$.

A rejeição de $H_0$ — ausência de efeito ARCH — direciona a série para o modelo ARIMA-GARCH, no qual o componente ARIMA captura a média condicional e o GARCH captura a dinâmica da variância residual \cite{bollerslev1986generalized}. Séries sem evidência de heterocedasticidade condicional avançam para a Etapa 4.

\subsection{Etapa 4 — Memória Longa}
\label{subsec:etapa4}

A quarta etapa avalia se a série apresenta dependência de longo alcance. O expoente de Hurst é estimado pelo método R/S, conforme a formulação original de \citeonline{hurst1951long-storage}. O resultado é confirmado pelo estimador GPH de \citeonline{geweke1983estimation}, que opera no domínio da frequência por meio de regressão log-periódograma e fornece estimativa do parâmetro de integração fracionária $d$.

Séries com $H > 0{,}7$ e parâmetro $d \in (0;\, 0{,}5)$ confirmado pelo GPH são atribuídas ao modelo ARFIMA. Conforme \citeonline{granger1980introduction}, a relação $H = d + 0{,}5$ conecta o expoente de Hurst ao parâmetro fracionário do ARFIMA. Séries que chegam à Etapa 4 sem indicação de memória longa são classificadas como lineares estacionárias e atribuídas ao modelo ARIMA.

% ----------------------------------------------------------
\section{Família de Modelos por Perfil}
\label{sec:familia-modelos}

A Tabela \ref{tab:familia-modelos} sintetiza os cinco perfis diagnósticos produzidos pelo protocolo e os modelos indicados para cada perfil.

\begin{table}[H]
  \centering
  \caption{Perfis diagnósticos e modelos indicados pelo protocolo.}
  \label{tab:familia-modelos}
  \begin{tabular}{p{3.5cm}p{6.5cm}p{3.5cm}}
    \toprule
    \textbf{Perfil} & \textbf{Critério de identificação} & \textbf{Modelo indicado} \\
    \midrule
    Linear estacionário  & ADF + KPSS concordam; BDS não rejeita; ARCH não rejeita; Hurst $\leq 0{,}7$ & ARIMA \\
    Sazonal              & Etapas 1–4 sem redirecionamento; $F_S > 0{,}64$ na Etapa 0 & SARIMA \\
    Não-linear           & BDS rejeita $H_0$ consistentemente                                          & LSTM \\
    Heterocedástico      & ARCH rejeita $H_0$                                                          & ARIMA-GARCH \\
    Memória longa        & $H > 0{,}7$ e GPH confirma $d \in (0;\, 0{,}5)$                            & ARFIMA \\
    \bottomrule
  \end{tabular}
  \fonte{Elaboração própria.}
\end{table}

A implementação de cada modelo segue as configurações descritas a seguir.

\textbf{ARIMA}: seleção de $(p, d, q)$ por minimização do AIC via busca em grade com $p, q \leq 5$, conforme a metodologia de \citeonline{box1970time}. O parâmetro $d$ é determinado pela Etapa 1 do protocolo.

\textbf{SARIMA}: inclui os componentes sazonais $(P, D, Q)_{12}$ adicionados ao ARIMA quando a Etapa 0 do protocolo detecta sazonalidade forte ($F_S > 0{,}64$) e as Etapas 1–4 não redirecionam a série para outro perfil. A seleção das ordens sazonais segue a metodologia de \citeonline{box1970time}.

\textbf{LSTM}: arquitetura de uma camada oculta com 64 unidades, janela de entrada de 24 passos, treinamento por até 50 épocas com \textit{early stopping} com paciência igual a 5, e otimizador Adam com taxa de aprendizado $10^{-3}$ \cite{hochreiter1997long}.

\textbf{ARIMA-GARCH}: o ARIMA captura a média condicional e o GARCH(1,1), conforme \citeonline{bollerslev1986generalized}, modela a variância condicional dos resíduos.

\textbf{ARFIMA}: estimação do parâmetro $d$ pelo método de máxima verossimilhança exata; ordens $p$ e $q$ selecionadas por AIC, conforme a formulação de \citeonline{granger1980introduction}.

% ----------------------------------------------------------
\section{Baselines de Comparação}
\label{sec:baseline}

Dois pipelines de referência são utilizados, ambos aplicados a todas as séries sem diagnóstico prévio.

\textbf{Baseline 1 — Auto-ARIMA.} Seleção automática de $(p, d, q)$ por minimização do AIC com teste de raiz unitária integrado, reproduzindo a prática comum em estudos de previsão automatizada \cite{hyndman2018forecasting}. Representa o paradigma de escolha por tentativa e erro sem triagem estatística prévia.

\textbf{Baseline 2 — Método Theta.} O Método Theta de \citeonline{assimakopoulos2000theta}, descrito na Seção \ref{subsec:theta}, é aplicado uniformemente a todas as séries. A escolha é motivada pelo desempenho histórico do método: vencedor da Competição M3 e entre os mais competitivos na M4, ele constitui o benchmark simples de referência mais relevante para o subconjunto mensal \cite{makridakis2018m4results}. A implementação utiliza a biblioteca \texttt{statsforecast} (Nixtla), que disponibiliza o método com tratamento automático de sazonalidade.

A comparação tripla — protocolo guiado, Auto-ARIMA e Theta — permite posicionar a contribuição em relação tanto à abordagem mais utilizada na prática quanto ao melhor método simples documentado na literatura M4.

% ----------------------------------------------------------
\section{Pipeline Experimental}
\label{sec:pipeline}

O experimento é executado nas seguintes etapas sequenciais:

\begin{enumerate}
  \item \textbf{Carregamento e particionamento}: as séries são carregadas do repositório oficial da Competição M4 \cite{makridakis2020m4}; a divisão treino/teste segue o protocolo original da competição.
  \item \textbf{Diagnóstico}: o protocolo de cinco etapas (Etapa 0 a Etapa 4) é aplicado individualmente a cada série da amostra de treino; cada série recebe um perfil diagnóstico e um modelo indicado.
  \item \textbf{Modelagem guiada}: o modelo indicado pelo protocolo é ajustado sobre a amostra de treino de cada série.
  \item \textbf{Modelagem com todos os candidatos}: os cinco modelos da família (ARIMA, SARIMA, LSTM, ARIMA-GARCH e ARFIMA) são ajustados sobre a amostra de treino de \emph{cada} série, independentemente do perfil diagnóstico. Esta etapa fornece os dados necessários para a análise de acurácia diagnóstica descrita na Seção \ref{sec:validacao-protocolo}.
  \item \textbf{Modelagem dos baselines}: Auto-ARIMA e Método Theta são ajustados sobre a amostra de treino de cada série.
  \item \textbf{Previsão}: todos os pipelines geram previsões para os 18 passos do horizonte de teste.
  \item \textbf{Avaliação}: sMAPE, MASE e OWA são calculados para cada série, conforme as Seções \ref{subsec:smape}, \ref{subsec:mase} e \ref{subsec:owa}, e agregados globalmente, por perfil diagnóstico, por domínio de origem e por faixa de comprimento da série, conforme Seção \ref{sec:analises-complementares}.
  \item \textbf{Análise estatística}: os resultados dos três pipelines (guiado, Auto-ARIMA, Theta) são comparados pelo teste de Wilcoxon para amostras pareadas \cite{wilcoxon1945individual}, com nível de significância $\alpha = 0{,}05$.
\end{enumerate}

% ----------------------------------------------------------
\section{Validação da Acurácia Diagnóstica}
\label{sec:validacao-protocolo}

A maioria dos trabalhos que propõem protocolos de diagnóstico verifica apenas se o método guiado produz menor erro agregado, sem examinar se os testes estatísticos de fato identificaram o modelo correto para cada série individualmente. Esta seção descreve uma análise secundária que responde a essa pergunta diretamente.

\subsection{Definição do Modelo Ótimo Empírico}
\label{subsec:oracle}

Para cada série $i$, define-se o \textbf{modelo ótimo empírico} (ou \textit{oracle}) como o modelo da família do protocolo que apresenta o menor OWA na amostra de teste:

\begin{equation}
  m^*_i = \arg\min_{m \,\in\, \mathcal{M}} \,\text{OWA}(m,\, i)
  \label{eq:oracle}
\end{equation}

onde $\mathcal{M} = \{\text{ARIMA}, \text{SARIMA}, \text{LSTM}, \text{ARIMA-GARCH}, \text{ARFIMA}\}$ é a família de modelos candidatos. O oracle não é um método aplicável na prática — ele requer conhecimento do futuro para ser calculado —, mas fornece o limite superior de desempenho alcançável pelo protocolo se o diagnóstico fosse perfeito.

\subsection{Taxa de Acerto e Arrependimento}
\label{subsec:taxa-acerto}

Seja $m^p_i$ o modelo atribuído pelo protocolo à série $i$. Definem-se duas métricas de acurácia diagnóstica:

\textbf{Taxa de acerto} ($\text{TA}$): fração de séries em que o protocolo identificou exatamente o modelo ótimo empírico:

\begin{equation}
  \text{TA} = \frac{1}{N} \sum_{i=1}^{N} \mathbf{1}\!\left[m^p_i = m^*_i\right]
  \label{eq:taxa-acerto}
\end{equation}

\textbf{Arrependimento médio} ($\bar{r}$): diferença média de OWA entre o modelo indicado pelo protocolo e o oracle, medindo o custo do erro de diagnóstico nas séries em que ele ocorre:

\begin{equation}
  \bar{r} = \frac{1}{N} \sum_{i=1}^{N} \left[\text{OWA}(m^p_i,\, i) - \text{OWA}(m^*_i,\, i)\right]
  \label{eq:arrependimento}
\end{equation}

Note que $\bar{r} \geq 0$ por construção; $\bar{r} = 0$ corresponderia a um protocolo de diagnóstico perfeito.

\subsection{Análise por Perfil}
\label{subsec:acerto-por-perfil}

As métricas $\text{TA}$ e $\bar{r}$ são calculadas globalmente e desagregadas por perfil diagnóstico, permitindo identificar em quais categorias o protocolo é mais e menos confiável. Uma baixa taxa de acerto em determinado perfil indica que o teste estatístico correspondente não discrimina bem o modelo ótimo para aquele tipo de série, o que fornece direção para refinamentos futuros do protocolo.

% ----------------------------------------------------------
\section{Análises Complementares}
\label{sec:analises-complementares}

Além da avaliação global descrita no Pipeline Experimental, três análises complementares são conduzidas para aprofundar a interpretação dos resultados e situar o protocolo no contexto da literatura da Competição M4.

\subsection{Desagregação por Domínio de Origem}
\label{subsec:analise-dominio}

O dataset M4 mensal rotula cada série com seu domínio de origem: Macro, Micro, Finanças, Indústria, Demografia e Outros \cite{makridakis2020m4}. O desempenho do protocolo é calculado separadamente para cada domínio, permitindo verificar se a vantagem do diagnóstico estatístico concentra-se em tipos específicos de série. A hipótese subjacente é que séries financeiras — reconhecidamente mais voláteis e não-lineares — devem beneficiar-se mais fortemente das Etapas 3 (ARCH/GARCH) e 2 (BDS/LSTM), enquanto séries demográficas, mais regulares e de longa duração, devem favorecer ARIMA e SARIMA. A desagregação por domínio permite verificar empiricamente essa expectativa e identificar possíveis pontos de falha do protocolo em domínios específicos.

\subsection{Estratificação por Comprimento da Série}
\label{subsec:analise-comprimento}

O comprimento da série de treino varia de 42 a 2.794 observações no subconjunto mensal. Séries curtas impõem restrições à potência dos testes estatísticos — particularmente ao teste BDS (que opera sobre resíduos de um AR ajustado) \cite{broock1996test} e ao estimador GPH (que utiliza regressão sobre o periodograma) \cite{geweke1983estimation} —, o que pode degradar a acurácia diagnóstica do protocolo em séries com menos observações. As séries são estratificadas em três faixas de comprimento — curta ($T < 100$), média ($100 \leq T < 500$) e longa ($T \geq 500$) —, e as métricas TA, $\bar{r}$ e OWA são calculadas em cada faixa. Esse recorte permite avaliar se o ganho do protocolo sobre os baselines é uniforme ou se depende de um volume mínimo de dados históricos.

\subsection{Posicionamento no Ranking M4}
\label{subsec:ranking-m4}

O OWA global obtido pelo protocolo é comparado com os resultados publicados dos 61 métodos que participaram oficialmente da Competição M4 \cite{makridakis2020m4}. A comparação não reproduz as condições exatas da competição — o protocolo opera sobre o subconjunto mensal, enquanto o ranking oficial cobre todas as frequências —, mas fornece uma referência qualitativa do nível de competitividade alcançado. O método vencedor da M4, ES-RNN de \citeonline{smyl2020hybrid}, e o segundo colocado, FFORMA de \citeonline{montero2020fforma}, servem como âncoras superiores da escala. O Método Theta e o Naive2, ambos disponíveis como benchmarks públicos, servem como âncoras inferiores. Resultados do protocolo posicionados entre o Theta e o FFORMA indicariam que o diagnóstico estatístico agrega valor real sem exigir o custo computacional de métodos híbridos complexos.

% ----------------------------------------------------------
\section{Ambiente Computacional}
\label{sec:ambiente}

Os experimentos são realizados em Python 3.11. As bibliotecas utilizadas são: \texttt{statsmodels} para os modelos ARIMA, ARFIMA e os testes ADF, KPSS e ARCH; \texttt{arch} para os modelos GARCH; \texttt{hurst} para estimação do expoente de Hurst e do parâmetro GPH; \texttt{pmdarima} para o Auto-ARIMA do baseline; \texttt{statsforecast} (Nixtla) para o Método Theta do baseline; e \texttt{tensorflow}/\texttt{keras} para o modelo LSTM. Os dados da Competição M4 são carregados via a biblioteca \texttt{M4py}. O código e os resultados são versionados em repositório de controle de versão.

% ============================================================
%  REFERÊNCIAS
% ============================================================
\bibliography{referencias}

\end{document}
